%%%%%%%%%%%%%%%%%%%%%%%%%%%%%%%%%%%%%%%%%%%%%%%%%%%
% Autor: Pedro Ivo da Cruz
% 
% Relatório Estudo Dirigido
% 
% Universidade Federal do Abc, Santo André, 2017
%
% A ideia deste template é simplificar a escrita do relatório de estudo 
% dirigido do programa de pós-graduação em Engenharia de Informação.
% Para isso, resolvi criar um arquivo de estilo, edinf.sty, que prove os 
% principais recursos de formatação de acordo com o guia de normalização
% acadêmica da UFABC, que pode ser encontrado no portal da biblioteca
% (www.portal.biblioteca.ufabc.edu.br).
%
% As instruções para o uso estão ao longo do documento.
% 
%%%%%%%%%%%%%%%%%%%%%%%%%%%%%%%%%%%%%%%%%%%%%%%%%%%

% \documentclass[a4paper, 12pt, twoside, openright]{report}
\documentclass[a4paper, 12pt]{report}

% O parâmetro twoside avisa o latex que o texto será impresso frente e verso. O parâmetro openright faz com que todo capítulo comece no anverso da página, ou seja, em uma página ímpar. Não há necessidade de usar este parâmetro no caso do relatório de estudo dirigido, mas está aí caso você deseje esta formatação.

%%%%%%%%%%%%%%%%%%%%%%%%%%%%%%%%%%%%%%%%%%%%%%%%%%%
% Pacotes necessários
%%%%%%%%%%%%%%%%%%%%%%%%%%%%%%%%%%%%%%%%%%%%%%%%%%%
\usepackage[english, brazil]{babel} % Português do Brasil, mas entende o inglês
\usepackage[utf8]{inputenc} % Codificação que permite acentos
\usepackage{edinf} % Estilo para relatório de estudo dirigido do programa de Engenharia da Informação da UFABC
\usepackage{amsmath}    % Equações
\usepackage{amsfonts}
\usepackage{amssymb}
\usepackage{amsthm} % Ambientes para exemplos, teoremas, provas e definições
\usepackage{mathtools}
\usepackage{array}      % Traz algumas funcionalidades úteis
\usepackage{verbatim}   % Traz algumas funcionalidades úteis
\usepackage{graphicx}   % Inclusão de figuras
\usepackage{epstopdf}   % Converte as imagens em EPS para PDF


%%%%%%%%%%%%%%%%%%%%%%%%%%%%%%%%%%%%%%%%%%%%%%%%%%%
% Seus pacotes aqui (opcionais)
%%%%%%%%%%%%%%%%%%%%%%%%%%%%%%%%%%%%%%%%%%%%%%%%%%%
\usepackage{cite}       % Melhorias nas citações
\usepackage{times}      % Usa fonte Times no texto
\usepackage[usenames, dvipsnames]{color} % Cores
\usepackage{multicol, multirow} % Pacote para auxiliar a montagem de tabelas
\usepackage{pdfpages}   % Facilita a inserção de páginas de PDFs externos
\usepackage{caption}    % Para importar o subcaption
\usepackage{subcaption} % Para usar subfiguras
\usepackage{algorithm}  % Ambiente para escrever algoritmos
\usepackage{algorithmic}
\usepackage{pgfgantt} % Gantt Charts

%%%%%%%%%%%%%%%%%%%%%%%%%%%%%%%%%%%%%%%%%%%%%%%%%%%
% Início do documento
%%%%%%%%%%%%%%%%%%%%%%%%%%%%%%%%%%%%%%%%%%%%%%%%%%%
\begin{document}

% Capa
% Acesse o arquivo e preencha os seus dados, e a data correta
\input{_1-capa}
\clearpage

%%%%%%%%%%%%%%%%%%%%%%%%%%%%%%%%%%%%%%%%%%%%%%%%%%%
% 				Elementos Pré-textuais
%%%%%%%%%%%%%%%%%%%%%%%%%%%%%%%%%%%%%%%%%%%%%%%%%%%
\pagenumbering{arabic} % Inicia a contagem das páginas em algarismos romanos.

% Folha de Rosto
% Acesse o arquivo e preencha os seus dados, e a data correta
\input{_2-folha_de_rosto}

% Resumo e Abstract
\input{_3-resumo}
\clearpage

% Lista de Figuras (Habilitar se quiser uma lista de figuras)
% \listoffigures
% \clearpage

% Sumário: a posição do sumário é relativa, alguns preferem antes das listas de figuras, tabelas, siglas, etc, outros gostam depois, logo antes dos elementos textuais. O guia de normalização recomenda colocá-lo aqui.
\tableofcontents
\thispagestyle{empty}
\clearpage

%%%%%%%%%%%%%%%%%%%%%%%%%%%%%%%%%%%%%%%%%%%%%%%%%%%
% 				Elementos Textuais
%%%%%%%%%%%%%%%%%%%%%%%%%%%%%%%%%%%%%%%%%%%%%%%%%%%
% \pagestyle{headings} % Use este se quiser cabeçalhos padrão do latex
\pagestyle{fancy} % Use este se quiser cabeçalho customizado pelo arquivo de estilo edinf.sty
% Defino este estilo aqui para apenas começar a mostrar a numeração das páginas aqui, de acordo com o guia de normalização.

% Introdução e Motivação
\input{1-introducao}
\clearpage

% Conclusão
\input{2-conclusao}
\clearpage

%%%%%%%%%%%%%%%%%%%%%%%%%%%%%%%%%%%%%%%%%%%%%%%%%%%
% 				Elementos Pós-textuais
%%%%%%%%%%%%%%%%%%%%%%%%%%%%%%%%%%%%%%%%%%%%%%%%%%%

% Bibliografia - Está por padrão a ABNT
\addcontentsline{toc}{chapter}{Bibliografia}
\renewcommand{\bibname}{Bibliografia}
\bibliography{referencias}

\end{document}
